\section{Comparación con los ajustes de constantes fundamentales}
\label{sec:comparacion}

La comparación externa evalúa las consecuencias de expresiones cuyos estados iniciales, regiones, coeficientes y acarreos han quedado determinados por la construcción anterior. Para cada referencia \(x_{\rm ref}\), con incertidumbre estándar \(u(x_{\rm ref})\), se informa el residuo normalizado
\[
 z=\frac{x_{\rm calc}-x_{\rm ref}}{u(x_{\rm ref})}.
\]
Esta razón expresa la diferencia respecto de cada centro tabulado en unidades de su incertidumbre estándar. Una interpretación probabilística requiere especificar el modelo estadístico; el contraste conjunto de los ajustes CODATA requiere además sus correlaciones y las condiciones de independencia que se utilicen.

Las expresiones utilizadas en esta comparación dan
\begin{align}
 \alpha&=0.007297352569283800997285105472380662999\ldots,\notag\\
 \frac{m_ec^2}{\mathrm{MeV}}&=0.510998950690000808608\ldots.
 \label{eq:valores-comparacion}
\end{align}
El segundo valor corresponde a la expresión electrónica corregida, con el factor completo de \(\Delta_4\) en la normalización basal y en el exponente de retorno. La composición del término basal con la memoria de retorno produce este valor antes de introducir la masa tabulada como referencia de contraste.

\subsection{Exactitud, incertidumbre y realización dimensional}

La comparación articula tres magnitudes con funciones distintas. La exactitud algebraica corresponde a las identidades demostradas en sus dominios. El error de evaluación numérica se controla mediante la precisión computacional. La incertidumbre experimental cuantifica el conocimiento de la magnitud medida según el ajuste utilizado y conserva ese significado al ampliar la evaluación decimal. Los dígitos adicionales de \eqref{eq:valores-comparacion} describen la expresión matemática; la información experimental está representada por los centros, incertidumbres y correlaciones del ajuste.

La realización dimensional del apartado~\ref{sec:electron} especifica la aplicación del carácter electrónico a energía y masa, con su escala y normalización. Esa aplicación determina la magnitud y las unidades sobre las que actúa el contraste. La construcción estructural establece el valor evaluado y sus relaciones; la comparación metrológica examina su correspondencia con la magnitud física así identificada. El orden causal conserva APP, TRIT y TPK como antecedentes de las expresiones evaluadas, seguidas por la realización dimensional y por los valores de referencia.

\subsection{CODATA 2018}

El ajuste 2018 contiene los valores centrales y las incertidumbres siguientes\cite{codata2018}:
\begin{center}
\begin{tabular}{@{}lcccr@{}}\toprule
Magnitud&Valor HMT&Referencia&\shortstack{Incertidumbre\\estándar}&\(z\)\\\midrule
\(\alpha\)&\(0.0072973525692838\ldots\)&\(0.0072973525693\)&\(1.1\times10^{-12}\)&\(-0.01473\)\\
\(m_ec^2/\mathrm{MeV}\)&\(0.5109989506900008\ldots\)&\(0.51099895000\)&\(1.5\times10^{-10}\)&\(+4.60001\)\\\bottomrule
\end{tabular}
\end{center}
La expresión de \(\alpha\) coincide con el valor central de 2018 al redondeo de las cifras mostradas. La diferencia, aproximadamente \(-1.62\times10^{-14}\), es unas \(0.0147\) veces la incertidumbre estándar impresa. La expresión electrónica corregida excede el centro de 2018 en aproximadamente \(6.90\times10^{-10}\,\mathrm{MeV}\), diferencia que persiste al redondeo indicado. Cada fila conserva así su concordancia decimal o su discrepancia cuantificada.

\subsubsection{Contraste de la ventana dodecafásica con el centro recíproco}
\label{sec:contraste-ventana-reciproco}

La ventana finita \(\alpha_{12}^{(0)}\) de
\eqref{eq:alpha-ventana-cero}, obtenida con \(c_{13}=0\), admite además
el contraste histórico con el recíproco del centro impreso de la
constante inversa. Definamos
\[
 R_{2018}:=137.035999084,\qquad
 \alpha_{\mathrm{rec},2018}:=R_{2018}^{-1}.
\]
En esta operación se evalúa el recíproco del racional decimal
\(R_{2018}\). Esta referencia se distingue del centro tabulado
\(0.0072973525693\) utilizado en la tabla anterior.

\begin{center}
\begin{tabular}{@{}ll@{}}
\toprule
Expresión & Evaluación aritmética\\
\midrule
\(\alpha_{12}^{(0)}\) &
\(0.007297352569283800997285105472380663\)\\
\(\alpha_{\mathrm{rec},2018}\) &
\(0.007297352569283800997285105472380663567972560701\ldots\)\\
\(\alpha_{12}^{(0)}-\alpha_{\mathrm{rec},2018}\) &
\(\simeq-5.679725607012965\times10^{-37}\)\\
\((\alpha_{12}^{(0)})^{-1}-R_{2018}\) &
\(\simeq+1.066588006664435\times10^{-32}\)\\
\(\bigl((\alpha_{12}^{(0)})^{-1}-R_{2018}\bigr)/R_{2018}\) &
\(\simeq+7.783268730800000\times10^{-35}\)\\
\bottomrule
\end{tabular}
\end{center}

La relación decimal precisa es
\[
 \alpha_{12}^{(0)}
 =10^{-36}\left\lfloor10^{36}\alpha_{\mathrm{rec},2018}\right\rfloor.
\]
Por tanto, la coincidencia histórica corresponde al truncamiento a
treinta y seis posiciones decimales del recíproco; su redondeo a esa
posición termina en \(664\), mientras la ventana termina en \(663\).
Las diferencias de la tabla cuantifican el contraste entre ambos
números. La sección infinita \(\alpha_A\) conserva su normalización
compatible y su frontera propias, distinguidas de esta ventana finita.

El ajuste expresa la constante inversa como
\(137.035999084(21)\), con incertidumbre estándar
\(u(R_{2018})=2.1\times10^{-8}\)\cite{codata2018}.
Los decimales obtenidos al invertir su centro impreso pertenecen a la
evaluación aritmética de ese racional; la incertidumbre experimental
conserva la indicada por el ajuste. La comparación precedente con el
centro tabulado de \(\alpha\) mantiene su referencia y su residuo
normalizado propios.


\subsection{CODATA 2022}

El ajuste 2022 modifica los valores centrales\cite{codata2022}:
\begin{center}
\begin{tabular}{@{}lcccr@{}}\toprule
Magnitud&Valor HMT&Referencia&\shortstack{Incertidumbre\\estándar}&\(z\)\\\midrule
\(\alpha\)&\(0.0072973525692838\ldots\)&\(0.0072973525643\)&\(1.1\times10^{-12}\)&\(+4.53073\)\\
\(m_ec^2/\mathrm{MeV}\)&\(0.5109989506900008\ldots\)&\(0.51099895069\)&\(1.6\times10^{-10}\)&\(+0.00000505\)\\\bottomrule
\end{tabular}
\end{center}
En esta edición la expresión electrónica coincide con el centro tabulado al redondeo, mientras que \(\alpha\) difiere aproximadamente \(4.53\) incertidumbres estándar impresas. La comparación de la exactitud física de los ajustes se fundamenta en sus datos experimentales, correlaciones e hipótesis. La fecha y la concordancia decimal con una expresión identifican propiedades del contraste; la evaluación de la proximidad al valor físico requiere ese análisis experimental completo.

\Needspace{29\baselineskip}
\subsection{Contraste de la constante de Planck reducida}
\label{sec:revision-planck-contraste}

El Sistema Internacional vigente desde el 20 de mayo de 2019 fija
exactamente \(h=6.62607015\times10^{-34}\,\mathrm{J\,s}\)
\cite{bipm2018}. Por tanto, su constante reducida es
\[
 \hbar_{\mathrm{SI}}=\frac{6.62607015}{2\pi}\,10^{-34}\,
 \mathrm{J\,s}
 =1.0545718176461563912\ldots\times10^{-34}\,\mathrm{J\,s}.
\]
Los puntos suspensivos corresponden a la expansión de una expresión
exacta, no a una incertidumbre experimental de \(h\) en el SI actual.
Este valor de referencia se utiliza aquí después de la evaluación de
los lectores de HMT; su función es comparativa.

Las dos coordenadas del retorno del artículo son
\begin{align*}
 H_5&=1.0545718176461544696\ldots,\\
 \eta_{\mathrm{ret}}&=1.0545718169150390611\ldots.
\end{align*}
En la identificación de \(\mathcal U_S\) con \(\mathrm{J\,s}\), la
diferencia relativa de la sección anterior respecto de
\(\hbar_{\mathrm{SI}}\) es aproximadamente \(-1.82\times10^{-15}\).
La de la sección posterior es aproximadamente \(-6.93\times10^{-10}\).
La segunda diferencia contiene el efecto del retorno que utiliza el
corrector electrónico; no equivale al error de evaluación de la
primera sección.

La precisión de estas comparaciones permite reconocer la proximidad
cuantitativa de la construcción de \(H_5\) y, simultáneamente, distinguir
proximidad de identidad exacta. En el presente manuscrito las expresiones
exhibidas no proporcionan una igualdad exacta con \(h/(2\pi)\).
El resultado demostrado sobre \(\Sigma_H\) es su década \(-34\);
el resultado analítico sobre \(H_5\) es la evaluación del carácter
declarado. La identificación física se examina en esta coordenatización y con las
diferencias explícitas que acabamos de dar. El valor SI no interviene
en la definición de la norma determinantal ni en la recurrencia regional.

